\section{Objectives}
\hspace{\parindent}
Building on these motivations, the primary objective of this dissertation is to study and improve the
robustness and effectiveness of centralized MRPP within the existing TEA* planning framework used
by the CRIIS and iiLab fleet coordination system at INESC TEC. This objective involves analyzing
how robot sequencing affects sequential TEA* planning and developing prioritization
heuristics capable of generating more informative robot sequences than random priority
selection.

Beyond the design of individual heuristics, this dissertation aims to evaluate how multiple priority
orderings can be explored within the same planning cycle through parallel execution. The proposed
work therefore seeks to reduce the dependence on a single robot sequence while preserving
compatibility with the existing coordination architecture.

Finally, the dissertation also aims to analyze temporal consistency issues observed during integrated
execution of the complete coordination stack, with particular attention to redundant replanning and
inconsistent processing of robot state updates. The anticipated outcome is a centralized
fleet coordination framework with improved empirical behavior in the evaluated planning scenarios, while
maintaining practical execution constraints.